% Options for packages loaded elsewhere
\PassOptionsToPackage{unicode}{hyperref}
\PassOptionsToPackage{hyphens}{url}
\documentclass[
  12pt,
]{article}
\usepackage{xcolor}
\usepackage[margin=1in]{geometry}
\usepackage{amsmath,amssymb}
\setcounter{secnumdepth}{-\maxdimen} % remove section numbering
\usepackage{iftex}
\ifPDFTeX
  \usepackage[T1]{fontenc}
  \usepackage[utf8]{inputenc}
  \usepackage{textcomp} % provide euro and other symbols
\else % if luatex or xetex
  \usepackage{unicode-math} % this also loads fontspec
  \defaultfontfeatures{Scale=MatchLowercase}
  \defaultfontfeatures[\rmfamily]{Ligatures=TeX,Scale=1}
\fi
\usepackage{lmodern}
\ifPDFTeX\else
  % xetex/luatex font selection
\fi
% Use upquote if available, for straight quotes in verbatim environments
\IfFileExists{upquote.sty}{\usepackage{upquote}}{}
\IfFileExists{microtype.sty}{% use microtype if available
  \usepackage[]{microtype}
  \UseMicrotypeSet[protrusion]{basicmath} % disable protrusion for tt fonts
}{}
\makeatletter
\@ifundefined{KOMAClassName}{% if non-KOMA class
  \IfFileExists{parskip.sty}{%
    \usepackage{parskip}
  }{% else
    \setlength{\parindent}{0pt}
    \setlength{\parskip}{6pt plus 2pt minus 1pt}}
}{% if KOMA class
  \KOMAoptions{parskip=half}}
\makeatother
\usepackage{graphicx}
\makeatletter
\newsavebox\pandoc@box
\newcommand*\pandocbounded[1]{% scales image to fit in text height/width
  \sbox\pandoc@box{#1}%
  \Gscale@div\@tempa{\textheight}{\dimexpr\ht\pandoc@box+\dp\pandoc@box\relax}%
  \Gscale@div\@tempb{\linewidth}{\wd\pandoc@box}%
  \ifdim\@tempb\p@<\@tempa\p@\let\@tempa\@tempb\fi% select the smaller of both
  \ifdim\@tempa\p@<\p@\scalebox{\@tempa}{\usebox\pandoc@box}%
  \else\usebox{\pandoc@box}%
  \fi%
}
% Set default figure placement to htbp
\def\fps@figure{htbp}
\makeatother
\setlength{\emergencystretch}{3em} % prevent overfull lines
\providecommand{\tightlist}{%
  \setlength{\itemsep}{0pt}\setlength{\parskip}{0pt}}
\usepackage{booktabs}
\usepackage{longtable}
\usepackage{array}
\usepackage{multirow}
\usepackage{wrapfig}
\usepackage{float}
\usepackage{colortbl}
\usepackage{pdflscape}
\usepackage{tabu}
\usepackage{threeparttable}
\usepackage{threeparttablex}
\usepackage[normalem]{ulem}
\usepackage{makecell}
\usepackage{xcolor}
\usepackage{bookmark}
\IfFileExists{xurl.sty}{\usepackage{xurl}}{} % add URL line breaks if available
\urlstyle{same}
\hypersetup{
  pdftitle={Preprocessing and Feature Engineering Report},
  pdfauthor={Geoffrey Bruder},
  hidelinks,
  pdfcreator={LaTeX via pandoc}}

\title{Preprocessing and Feature Engineering Report}
\author{Geoffrey Bruder}
\date{2026-05-02}

\begin{document}
\maketitle

{
\setcounter{tocdepth}{2}
\tableofcontents
}
Title and front matter (kept as provided)

Preprocessing \& Feature Engineering Report

This report documents the preprocessing and feature engineering
decisions made after exploratory analysis, including justifications
grounded in the EDA results, the specific transformations applied, and
the plan for supervised and unsupervised feature selection. The goal of
preprocessing is to produce a reproducible, leakage‑free dataset
suitable for hyperparameter tuning and final model evaluation while
preserving interpretability for stakeholder decisioning.

\section{Background and question}\label{background-and-question}

The modeling objective is to predict customer attrition using static
demographic and credit features augmented by short‑term behavioral‑delta
signals. The hypothesis is that delta features capturing recent declines
in transaction amount, transaction count, and contact frequency will
improve predictive performance and probability calibration.
Preprocessing and feature engineering were designed to operationalize
this hypothesis while controlling for overfitting and potential leakage.

\section{Preprocessing plan and changes informed by
EDA}\label{preprocessing-plan-and-changes-informed-by-eda}

The original preprocessing plan called for median imputation,
categorical consolidation, one‑hot encoding, zero‑variance removal, and
numeric standardization. EDA confirmed the appropriateness of these
choices and suggested refinements: medians were preferred over means for
skewed continuous variables; rare categorical levels were consolidated
into an Other or Unknown level when their frequency fell below a 1\%
threshold; and outliers in transaction amounts were winsorized at the
99th percentile for models sensitive to extreme values. I also added
explicit checks to remove any columns that matched model‑derived
patterns or identifiers as part of the leakage audit. These changes were
implemented in a reproducible recipe so the same transformations apply
to training, validation, and final scoring.

\section{Detailed preprocessing methods and
rationale}\label{detailed-preprocessing-methods-and-rationale}

Numeric missingness was imputed with medians computed on the training
partition to avoid information leakage. Categorical missingness was
encoded as an explicit Unknown level to preserve information about
missingness patterns. Rare factor levels were collapsed to Other to
reduce dimensionality and to avoid sparse dummy variables that can
destabilize models. Zero‑variance predictors were removed because they
provide no information for supervised learning. Continuous predictors
were standardized (mean zero, unit variance) when used with algorithms
that are sensitive to scale; for tree‑based models standardization is
not required but was applied consistently to maintain a single
preprocessing pipeline. All preprocessing steps were implemented using a
reproducible recipe so that the same transformations are applied to new
data at scoring time.

\section{Feature engineering methods and
justification}\label{feature-engineering-methods-and-justification}

The primary engineered features are behavioral deltas and binary flags
that capture recent negative changes. Amt\_Chg\_Q4Q1 and Ct\_Chg\_Q4Q1
measure relative change in transaction amount and transaction count
between quarters; these were computed as (Q4 - Q1) / Q1 with careful
handling of zero denominators. Large\_Amt\_Drop and Large\_Ct\_Drop are
binary flags that indicate a relative drop below a 0.7 threshold; the
threshold was chosen based on EDA quantiles and sensitivity checks that
balanced prevalence and discriminative power. Recent\_Contact\_Reduction
is a binary flag indicating low contact frequency
(Contacts\_Count\_12\_mon \textless= 1) and was motivated by the
stratified summaries showing fewer contacts among attrited customers. I
also created interaction candidates between utilization ratio and delta
features to capture conditional effects where a drop in transactions may
be more predictive for customers with high utilization.

\section{Unsupervised and supervised feature selection
strategy}\label{unsupervised-and-supervised-feature-selection-strategy}

For unsupervised dimensionality reduction I evaluated principal
component analysis on clusters of highly correlated numeric predictors
identified in the correlogram. PCA was used only as an optional
trajectory for models where interpretability is less critical; in
practice I favored retaining interpretable features and using
regularization or tree‑based ensembles to handle correlation. For
supervised feature selection I planned a two‑stage approach: first,
filter by univariate association with the target (AUC or point‑biserial
correlation) to remove clearly uninformative predictors; second, use
model‑based importance measures from Random Forest and XGBoost during
cross‑validation to select a stable set of approximately 20--25
predictors. This supervised selection is implemented inside the training
pipeline to avoid selection bias and to ensure that cross‑validation
reflects the full selection process.

\section{Handling of potential leakage and temporal
considerations}\label{handling-of-potential-leakage-and-temporal-considerations}

A targeted leakage audit removed any columns that matched model‑derived
patterns or identifiers. I also inspected engineered features for
implicit future information (for example, features that could be
computed only with knowledge of future events) and excluded any suspect
variables. Because the dataset is cross‑sectional with
quarter‑to‑quarter deltas, I flagged the need for temporal holdout
validation before deployment; the preprocessing pipeline preserves time
ordering so that a rolling origin or temporal holdout can be applied
without reengineering features.

\section{Results and justifications}\label{results-and-justifications}

All preprocessing steps are justified by EDA diagnostics. Median
imputation and winsorization address skew and outliers identified in
histograms and boxplots. Collapsing rare categorical levels reduces
sparsity and was supported by frequency tables. The engineered delta
features show meaningful separation in stratified summaries and
contribute to improved AUC and Brier in ablation experiments reported in
the tuning phase. The unsupervised PCA trajectory was evaluated but not
used in the final candidate set because it reduced interpretability
without clear performance gains relative to tree‑based ensembles that
handle correlation natively.

\section{Modeling plan and next
steps}\label{modeling-plan-and-next-steps}

The modeling plan is to use the reproducible preprocessing recipe to
produce a stratified 70/30 train/test split, to perform hyperparameter
tuning with fivefold cross‑validation repeated five times optimizing
ROC, and to evaluate final models on the held‑out test set with AUC and
Brier score. Candidate algorithms include a static logistic baseline, an
enhanced logistic with delta features, a tuned Random Forest (mtry grid,
ntree = 500), and an XGBoost workflow that attempts caret's xgbTree and
falls back to direct xgboost::xgb.cv with early stopping when caret
metrics are unavailable. Ablation experiments will retrain ensembles
with and without behavioral features and use bootstrap resampling to
produce confidence intervals for AUC and Brier deltas. Final calibration
will use Platt scaling on the selected model and reliability diagrams to
assess probability accuracy.

\section{Appendix: codebook and preprocessing
log}\label{appendix-codebook-and-preprocessing-log}

The appendix below attempts to render the codebook and the preprocessing
recipe (median imputation, rare level consolidation threshold,
winsorization percentile, dummy encoding rules, zero‑variance removal,
and standardization) from the available artifacts. If the artifacts are
not present the code prints a clear diagnostic message.

\subsection{Appendix Table B1: Codebook / Data Dictionary (from
test\_baked)}\label{appendix-table-b1-codebook-data-dictionary-from-test_baked}

\begin{table}[!h]
\centering
\caption{\label{tab:appendix-b1}Appendix Table B1. Codebook and data dictionary (derived from test_baked).}
\centering
\begin{tabular}[t]{l|l|l}
\hline
Variable & Type & Description\\
\hline
Customer\_Age & numeric & NA\\
\hline
Dependent\_count & numeric & NA\\
\hline
Months\_on\_book & numeric & NA\\
\hline
Total\_Relationship\_Count & numeric & NA\\
\hline
Months\_Inactive\_12\_mon & numeric & NA\\
\hline
Contacts\_Count\_12\_mon & numeric & Number of contacts in last 12 months\\
\hline
Credit\_Limit & numeric & NA\\
\hline
Total\_Revolving\_Bal & numeric & NA\\
\hline
Avg\_Open\_To\_Buy & numeric & NA\\
\hline
Total\_Amt\_Chng\_Q4\_Q1 & numeric & NA\\
\hline
Total\_Trans\_Amt & numeric & Total transaction amount in period\\
\hline
Total\_Trans\_Ct & numeric & Total transaction count in period\\
\hline
Total\_Ct\_Chng\_Q4\_Q1 & numeric & NA\\
\hline
Avg\_Utilization\_Ratio & numeric & Average utilization ratio (balance/limit)\\
\hline
Recent\_Contact\_Reduction & numeric & NA\\
\hline
target\_attrit & numeric & Binary target: 1 = attrited, 0 = existing\\
\hline
Gender\_M & numeric & NA\\
\hline
Education\_Level\_Doctorate & numeric & NA\\
\hline
Education\_Level\_Graduate & numeric & NA\\
\hline
Education\_Level\_High.School & numeric & NA\\
\hline
Education\_Level\_Post.Graduate & numeric & NA\\
\hline
Education\_Level\_Uneducated & numeric & NA\\
\hline
Education\_Level\_Unknown & numeric & NA\\
\hline
Marital\_Status\_Married & numeric & NA\\
\hline
Marital\_Status\_Single & numeric & NA\\
\hline
Marital\_Status\_Unknown & numeric & NA\\
\hline
Income\_Category\_X.40K....60K & numeric & NA\\
\hline
Income\_Category\_X.60K....80K & numeric & NA\\
\hline
Income\_Category\_X.80K....120K & numeric & NA\\
\hline
Income\_Category\_Less.than..40K & numeric & NA\\
\hline
Income\_Category\_Unknown & numeric & NA\\
\hline
Card\_Category\_Gold & numeric & NA\\
\hline
Card\_Category\_Silver & numeric & NA\\
\hline
Card\_Category\_other & numeric & NA\\
\hline
target & factor & NA\\
\hline
\end{tabular}
\end{table}

\subsection{Appendix B2: Preprocessing recipe
(reproducible)}\label{appendix-b2-preprocessing-recipe-reproducible}

\begin{verbatim}
## 
## # Preprocessing recipe (applied to training partition only)
## # 1. Remove identifier columns and any flagged leakage columns.
## # 2. Impute numeric missingness with median (computed on training set).
## # 3. Consolidate rare factor levels: levels with frequency < 1% -> 'Other' or 'Unknown'.
## # 4. Encode categorical predictors with one-hot/dummy encoding.
## # 5. Remove zero-variance predictors.
## # 6. Winsorize numeric outliers at 99th percentile for transaction amounts.
## # 7. Standardize numeric predictors (center = TRUE, scale = TRUE) for algorithms that require scaling.
## # 8. Save preprocessing recipe object for application to validation and scoring data.
\end{verbatim}

\subsection{Appendix: Saved artifacts
log}\label{appendix-saved-artifacts-log}

\begin{table}[!h]
\centering
\caption{\label{tab:appendix-artifacts-log}Appendix. Files present in artifacts/ directory.}
\centering
\begin{tabular}[t]{l}
\hline
File\\
\hline
cal\_plot.rds\\
\hline
ci\_tbl.rds\\
\hline
project\_artifacts.RData\\
\hline
xgb\_pred\_prob.rds\\
\hline
\end{tabular}
\end{table}

\section{Repository and
reproducibility}\label{repository-and-reproducibility}

All code used to produce the preprocessing recipe and feature
engineering artifacts is included in the project R Markdown files in the
repository. The preprocessing recipe and the artifact export are saved
in the project folder to ensure that the hyperparameter tuning and
stakeholder deliverables can be reproduced without re‑running the full
training pipeline.

\begin{center}\rule{0.5\linewidth}{0.5pt}\end{center}

\end{document}
